Confira cada grupo mesmo quando parecer correto. Reproduza a falha e reduza o caso; compare estados intermediários com uma simulação manual ou força bruta.
\begin{itemize}[leftmargin=*,label={}]
\item \ifextended\else\begin{minipage}[t]{\linewidth}\fi\textbf{Enunciado e modelo}
\begin{itemize}[leftmargin=*,label={\(\square\)}]
\item Releia o que deve ser contado/otimizado: ordem importa? Repetições são distintas? A escolha vazia é permitida?
\item Confira as hipóteses do algoritmo: pesos negativos, monotonicidade, conectividade, empates. Elas vêm do enunciado ou foram presumidas?
\item Siga à mão um caso pequeno; confira o significado de cada estado, transição, base e resposta extraída.
\end{itemize}
\ifextended\else\end{minipage}\fi

\item \ifextended\else\begin{minipage}[t]{\linewidth}\fi\textbf{Limites e índices}
\begin{itemize}[leftmargin=*,label={\(\square\)}]
\item Force menor tamanho permitido, um elemento, todos iguais, valores extremos e resposta inexistente.
\item Confira cada acesso: 0/1-based, $n/m$ trocados, $i-1$, $i+1$, sentinelas $-1/n$, tamanho $n+1/2n$.
\item Refaça os extremos dos intervalos: $[l,r)$ ou $[l,r]$? Teste vazio, unitário e intervalo inteiro; confira $<$ versus $\le$.
\end{itemize}
\ifextended\else\end{minipage}\fi

\item \ifextended\else\begin{minipage}[t]{\linewidth}\fi\textbf{Aritmética}
\begin{itemize}[leftmargin=*,label={\(\square\)}]
\item Verifique o tipo de cada produto, soma e shift \emph{antes} da atribuição; cabe no tipo? Alguma operação usa \texttt{INF}?
\item Confira módulo, resíduos negativos e divisão modular (o inverso existe?). Na divisão inteira, o arredondamento é o desejado?
\item Teste denominador zero, raiz de negativo, \texttt{NaN}/infinito e cancelamento numérico; confira tolerância e precisão da saída.
\end{itemize}
\ifextended\else\end{minipage}\fi

\item \ifextended\else\begin{minipage}[t]{\linewidth}\fi\textbf{Estado e atualizações}
\begin{itemize}[leftmargin=*,label={\(\square\)}]
\item Rode o mesmo caso sozinho e após outro: globais, containers, flags e lazy tags foram reinicializados?
\item Cada valor foi inicializado antes de ser lido? Algum estado impossível foi usado como válido?
\item A atualização usa o estado anterior ou já alterado? Confira direção do loop, ordem de composição e rollback; referências/iteradores continuam válidos?
\end{itemize}
\ifextended\else\end{minipage}\fi

\item \ifextended\else\begin{minipage}[t]{\linewidth}\fi\textbf{Representação e ordem}
\begin{itemize}[leftmargin=*,label={\(\square\)}]
\item Teste componentes desconexas, isolados, self-loops e arestas paralelas; direção, aresta reversa e identificação pelo pai estão corretas?
\item Force empates: a ordenação tem desempate correto e comparador estrito? Eventos na mesma coordenada entram/saem na ordem certa?
\item Após ordenar/comprimir, os IDs e pesos ainda correspondem? A compressão perdeu distâncias necessárias? A reconstrução produz uma solução válida?
\end{itemize}
\ifextended\else\end{minipage}\fi

\item \ifextended\else\begin{minipage}[t]{\linewidth}\fi\textbf{Tempo e memória}
\begin{itemize}[leftmargin=*,label={\(\square\)}]
\item Cada loop progride? Recursão revisita estados ou excede a pilha? Teste cadeia, estrela e entrada adversária para a estrutura.
\item Conte cópias, alocações, limpeza por teste e fatores extras de $\log n$; confira o custo total de todos os casos e o argumento amortizado.
\item Calcule o tamanho real dos buffers/estados; confira se containers crescem sem limite ou se estruturas locais grandes ocupam a pilha.
\end{itemize}
\ifextended\else\end{minipage}\fi

\item \ifextended\else\begin{minipage}[t]{\linewidth}\fi\textbf{Execução e saída}
\begin{itemize}[leftmargin=*,label={\(\square\)}]
\item Recompile o arquivo correto e rode com sanitizers. Leia todos os campos e casos? Removeu prints de debug? Formato e precisão estão corretos?
\item Em stress testing, o gerador respeita as restrições e a força bruta é independente? Respostas múltiplas exigem validar a solução, não igualdade textual.
\end{itemize}
\ifextended\else\end{minipage}\fi
\end{itemize}
